\documentclass[../main.tex]{subfiles}
\graphicspath{{\subfix{../images/}}}

\begin{document}

\section{TODO}
\begin{itemize}
\item reverse control/cases in table 1
\end{itemize}

\section{Introduction}
\begin{itemize}
    \item climate change $\rightarrow$ increasing frequency of PSPS events and wildfire smoke
    \item relationship between power outages and health
    \item relationship between wf and health 
    \item concurrent events and health -- introduce psps events bc it is a manifestation of this direct link between wfs and outages. 
    \item study objectives
    \begin{itemize}
        \item examine the association between PSPS events and health outcomes.
        \item assess interaction effects with wildfire smoke on health outcomes.
    \end{itemize}
\end{itemize}

\section{Methods}
    \subsection{data}
    \begin{itemize}
        \item exposure 
        \begin{itemize}
          \item PSPS data description (description of psps events, description of data source, description of categories, specification of spatiotemporal res)
          \item wildfire smoke data description
        \end{itemize}
        \item outcome
        \begin{itemize}
          \item HCAI data data description (hosp admissions + ED visits, adults, CA)
        \end{itemize}
    \end{itemize}

    \subsection{study population}
    \begin{itemize}
        \item CA adults
    \end{itemize}

    \subsection{study design}
    \begin{itemize}
        \item case-crossover design
    \end{itemize}
    
    \subsection{statistical analysis}
    \begin{itemize}
        \item conditional logistic regression 
        \item interaction term between PSPS severity and wildfire smoke exposure.
    \end{itemize}


\section{Results}
\begin{itemize}
    \item overview (report main effects first, then co-occurrence, for each)
    \begin{itemize}
        \item description of study population (e.g., table 1).
        \item interactions of severe PSPS events and wildfire smoke for respiratory morbidity (1.26, 95\% CI 1.03 - 1.53)
        \item when analyzing copd alone, strong effects for just wf or just psps, stronger than resp overall. but dont detect an effect for co-occurrence of thse and its unclear whether we are underpowered or if the relationship doesnt exist here. 
        \item cardio para -- report null findings 
    \end{itemize}
    \item tables/figures
    \begin{itemize}
        \item \textbf{table 1}: Descriptive statistics of hospital admissions and ED visits by PSPS severity and wildfire smoke exposure.
        \item \textbf{table 2}: descriptives of wf exp and psps exp overall and their co-occurrence (avg wf exposure during psps events and what is the range) 
        \item \textbf{figure 1}: map of PSPS and avg wf smoke pm during psps events -- play with this, maybe only show psps  
        \item \textbf{figure 2}: plots showing effect estimates (box plots? could we do something more creative?) 
    \end{itemize}
\end{itemize}

\section{Discussion}
\begin{itemize}
    \item summary of findings 
    \begin{itemize}
        \item significant interaction effects for respiratory morbidity and their public health implications.
        \item Non-significant effects for COPD, cardiovascular, and psychiatric morbidity.
        \item TBD
    \end{itemize}
    \item limitations
    \item policy implications and next steps 
    \item call for future research on co-occurrence of weather and psps events 
    
\end{itemize}

\clearpage
\setcounter{section}{0}
